% ============================================================
% Cheat sheet, argomento 2: il diagramma delle classi, le relazioni (dispensa, capitolo 7)
% ============================================================
\argomento{Il diagramma delle classi: le relazioni}{Dispensa: capitolo 7}[Le relazioni fra classi]

\begin{guida}{Le linee e le frecce in un colpo d'occhio}
\centering
\renewcommand{\arraystretch}{1.45}
\begin{tabular}{@{}c >{\bfseries}l l l@{}}
  \textbf{Disegno} & \textbf{Relazione} & \textbf{Si legge} & \textbf{In Java} \\ \hline
  \tikz[baseline=-0.6ex]{\draw[assocl] (0,0) -- (1.9,0);} & associazione & A è collegata a B & un campo (in uno o in entrambi i versi) \\
  \tikz[baseline=-0.6ex]{\draw[navig] (0,0) -- (1.9,0);} & associazione navigabile & da A si arriva a B & A ha un campo di tipo B \\
  \tikz[baseline=-0.6ex]{\draw[aggr] (0,0) -- (1.9,0);} & aggregazione & A (rombo) è fatta di B & campo o collezione di B, parti condivisibili \\
  \tikz[baseline=-0.6ex]{\draw[compos] (0,0) -- (1.9,0);} & composizione & B non esiste senza A & B creata e posseduta da A \\
  \tikz[baseline=-0.6ex]{\draw[gen] (0,0) -- (1.9,0);} & generalizzazione & A è un B & \code{class A extends B} \\
  \tikz[baseline=-0.6ex]{\draw[realizza] (0,0) -- (1.9,0);} & realizzazione & A implementa B & \code{class A implements B} \\
  \tikz[baseline=-0.6ex]{\draw[dep] (0,0) -- (1.9,0);} & dipendenza & A usa B & parametro, variabile locale, \code{new B()}, \code{B.m()} \\
  \tikz[baseline=-0.6ex]{\draw[tratt] (0,0) -- (1.9,0);} & linea tratteggiata & nota o classe associativa & --- \\
\end{tabular}

\smallskip
\raggedright La punta \textbf{vuota a triangolo} vuol dire sempre «è un tipo di» (linea continua: eredita, tratteggiata:
implementa). Il \textbf{rombo} sta sempre dal lato del \textbf{tutto}. Il \textbf{tratteggio} con la punta aperta è
sempre una \textbf{dipendenza}.
\end{guida}

\voce{associazione}{Associazione}{linea con un nome e il verso di lettura}
  {\begin{tikzpicture}
     \node[classe1=1.5cm] (p) at (0,0) {Persona};
     \node[classe1=1.7cm] (a) at (4.2,0) {Automobile};
     \draw[assocl] (p) -- node[above, font=\scriptsize]{Possiede $\blacktriangleright$} (a);
   \end{tikzpicture}}
  {Una relazione fra le istanze di due classi. Il triangolino dice solo come si \textbf{legge} il nome: «la Persona
   possiede l'Automobile». Non dice chi conosce chi (quello è il verso di navigazione, \vv{navigazione}).}

\voce{ruolo}{Ruolo}{il nome di un estremo}
  {\begin{tikzpicture}
     \node[classe1=1.5cm] (p) at (0,0) {Persona};
     \node[classe1=1.7cm] (a) at (4.2,0) {Automobile};
     \draw[assocl] (p) -- (a);
     \node[ruolo, anchor=south west] at ($(p.east)+(0.05,0.02)$) {+proprietario};
     \node[ruolo, anchor=north east] at ($(a.west)+(-0.05,-0.02)$) {+automobili};
   \end{tikzpicture}}
  {Il ruolo si scrive vicino alla classe che lo ricopre e diventa il \textbf{nome del campo} dall'altra parte:
   \code{Persona proprietario} sta in \code{Automobile}, \code{List<Automobile> automobili} in \code{Persona}.}

\voce{molteplicitaassoc}{Molteplicità di un'associazione}{si scrive vicino alla classe che conta}
  {\begin{tikzpicture}
     \node[classe1=1.5cm] (p) at (0,0) {Persona};
     \node[classe1=1.7cm] (a) at (4.2,0) {Automobile};
     \draw[assocl] (p) -- node[above, font=\scriptsize]{Possiede $\blacktriangleright$}
        node[molt, below, pos=0.08]{1} node[molt, below, pos=0.9]{1..*} (a);
   \end{tikzpicture}}
  {Il \code{1..*} vicino ad Automobile si legge: «\textit{una} Persona possiede \textit{da 1 a molte}
   Automobili». Il \code{1} vicino a Persona: «\textit{un'}Automobile ha \textit{un solo} proprietario». Uno a molti.
   I costruttori devono rispettarle (\vv{unomolti}).}

\voce{navigazione}{Verso di navigazione}{freccia aperta; croce = non navigabile}
  {\begin{tikzpicture}
     \node[classe1=1.5cm] (p) at (0,0) {Persona};
     \node[classe1=1.7cm] (a) at (4.2,0) {Automobile};
     \draw[navig] (p) -- (a);
     \estO{(a.west)}{0..1}{auto}
     \noNav{($(p.east)+(0.22,0)$)}{0}
   \end{tikzpicture}}
  {Da una Persona si trova la sua automobile; da un'Automobile \textbf{non} si risale alla persona (la croce). In
   Java: \code{Automobile auto} sta in \code{Persona}, e \code{Automobile} non ha campi di tipo \code{Persona}. Senza
   frecce la navigazione non è specificata, oppure va nei due versi. È una scelta di progetto: nei diagrammi
   concettuali non si mette.}

\voce{aggregazione}{Aggregazione}{rombo vuoto dal lato del tutto}
  {\begin{tikzpicture}
     \node[classe1=1.5cm] (v) at (0,0) {Veicolo};
     \node[classe1=1.5cm] (r) at (4.0,0) {Ruota};
     \draw[aggr] (v) -- node[molt, above, pos=0.9]{*} node[molt, below, pos=0.22]{0..1} (r);
   \end{tikzpicture}}
  {Relazione \textbf{tutto--parti}: le parti possono esistere anche da sole (una ruota di scorta) e passare da un tutto
   all'altro. Se dal lato del tutto non c'è niente, vale \code{0..1}. Si implementa come un'associazione
   (\vv{aggregazionecodice}).}

\voce{composizione}{Composizione}{rombo pieno: le parti muoiono con il tutto}
  {\begin{tikzpicture}
     \node[classe1=1.5cm] (c) at (0,0) {Casa};
     \node[classe1=1.5cm] (s) at (4.0,0) {Stanza};
     \draw[compos] (c) -- node[molt, above, pos=0.9]{*} node[molt, below, pos=0.22]{1} (s);
   \end{tikzpicture}}
  {Aggregazione \textbf{forte}: una parte appartiene a \textbf{un solo} tutto e non esiste senza di lui. Sottintesa la
   molteplicità \code{1} dal lato del tutto. Nel codice il tutto crea le parti e non le cede
   (\vv{composizionecodice}).}

\voce{generalizzazione}{Generalizzazione}{triangolo vuoto verso la superclasse}
  {\begin{tikzpicture}
     \node[classe1=1.6cm] (a) at (0,0) {\textit{Animale}};
     \node[classe1=1.6cm] (c) at (-1.3,-1.6) {Cane};
     \node[classe1=1.6cm] (g) at (1.3,-1.6) {Gatto};
     \coordinate (t) at (0,-0.85);
     \draw[gen] (t) -- (a.south);
     \draw[assocl] (c.north) |- (t) -| (g.north);
     \node[ruolo, anchor=west] at (0.1,-0.65) {specie};
     \node[ruolo, anchor=west] at (2.25,-1.1) {\{incomplete\}};
   \end{tikzpicture}}
  {«Cane \textbf{è un} Animale»: \code{class Cane extends Animale}. La sottoclasse eredita attributi e operazioni.
   Il \textbf{discriminatore} (\code{specie}) dice il criterio; \code{\{incomplete\}} (o una linea tratteggiata che
   taglia le frecce) dice che esistono altre sottoclassi non disegnate.}

\voce{realizzazione}{Realizzazione}{tratteggiata con il triangolo vuoto}
  {\begin{tikzpicture}
     \umlclass[3.3cm]{i}{(0,0)}{\stereotipo{interface}{Cassiere}}{}{+ preleva(importo: int)}
     \node[classe1=1.5cm] (e) at (-1.4,-2.0) {Impiegato};
     \node[classe1=1.5cm] (b) at (1.4,-2.0) {Bancomat};
     \draw[realizza] (e.north) |- (0,-1.3) -- (i.south);
     \draw[realizza] (b.north) |- (0,-1.3);
   \end{tikzpicture}}
  {\code{class Bancomat implements Cassiere}: la classe fornisce il codice di tutte le operazioni dell'interfaccia.
   Classi di gerarchie diverse possono realizzare la stessa interfaccia, e chi usa un \code{Cassiere} non sa quale
   classe c'è dietro.}

\voce{lollipop}{Pallina e presa}{interfaccia fornita e richiesta}
  {\begin{tikzpicture}
     \node[classe1=1.9cm] (m) at (0,0) {MotoreRicerca};
     \node[classe1=1.6cm] (b) at (4.1,0) {Browser};
     \draw[assocl] (m.east) -- (2.15,0);
     \palla{(2.26,0)}
     \presa{(2.26,0)}{180}
     \draw[assocl] (2.06,0) -- (b.west);
     \node[ruolo] at (2.26,-0.42) {Cerca};
   \end{tikzpicture}}
  {La \textbf{pallina} (\textit{lollipop}) è l'interfaccia che \code{MotoreRicerca} \textbf{fornisce} (la implementa); la
   \textbf{presa} (\textit{socket}) è quella che \code{Browser} \textbf{richiede} (la usa: è una dipendenza). Equivale a
   una realizzazione più una dipendenza verso \code{«interface» Cerca}.}

\voce{dipendenza}{Dipendenza}{freccia tratteggiata, con uno stereotipo}
  {\begin{tikzpicture}
     \node[classe1=1.2cm] (a) at (0,0) {A};
     \node[classe1=1.2cm] (b) at (3.6,0.55) {B};
     \node[classe1=1.2cm] (c) at (3.6,-0.55) {C};
     \draw[dep] (a.east) -- node[stereo, above, sloped]{«call» «use»} (b.west);
     \draw[dep] (a.east) -- node[stereo, below, sloped]{«create»} (c.west);
   \end{tikzpicture}}
  {A \textbf{usa} B durante l'esecuzione, senza tenerne un riferimento in un campo: \textbf{«call»} chiama un suo
   metodo, \textbf{«use»} ne usa un attributo, \textbf{«create»} ne crea un'istanza (\code{new}). Se B cambia, A può non
   compilare più: è un accoppiamento. Si vede nel progetto di dettaglio (\vv{dipendenzecodice}).}

\voce{classeassociativa}{Classe associativa}{collegata all'associazione con il tratteggio}
  {\begin{tikzpicture}
     \node[classe1=1.5cm] (s) at (0,0) {Studente};
     \node[classe1=1.5cm] (e) at (4.0,0) {Esame};
     \draw[assocl] (s) -- node[molt, above, pos=0.1]{*} node[molt, above, pos=0.9]{*} (e);
     \umlclass[1.5cm]{v}{(2.0,-1.25)}{Voto}{valore: int}{}
     \draw[tratt] (2.0,0) -- (v.north);
   \end{tikzpicture}}
  {Un attributo (il voto) che riguarda la \textbf{coppia} studente--esame e non si può dare a nessuno dei due. Tipica
   dei molti a molti. Equivale a una classe normale con due associazioni uno a molti
   (\code{Studente 1 -- * Voto * -- 1 Esame}), che però ammette più voti per la stessa coppia (\vv{associativacodice}).}

\voce{riflessiva}{Associazione riflessiva}{una classe collegata a se stessa}
  {\begin{tikzpicture}
     \node[classe1=1.8cm, minimum height=0.8cm] (c) at (0,0) {Corso};
     \draw[assocl] (0.9,0.15) -- (1.6,0.15) -- (1.6,1.0) -- (0.45,1.0) -- (0.45,0.4);
     \node[molt] at (1.2,0.3) {*};
     \node[molt] at (0.3,0.6) {*};
     \node[ruolo, anchor=west] at (1.65,0.85) {propedeutico a $\blacktriangleright$};
     \node[ruolo, anchor=west] at (1.65,0.3) {\code{List<Corso>}};
   \end{tikzpicture}}
  {Un \code{Corso} è propedeutico ad altri corsi: in Java un campo il cui tipo è la classe stessa
   (\code{List<Corso> propedeutici}, \code{Persona coniuge}). Asimmetrica (propedeutico a) o simmetrica (in
   alternativa a). I ruoli distinguono i due estremi (\vv{riflessivacodice}).}

\voce{qualificata}{Associazione qualificata}{la chiave per trovare l'oggetto \oltre}
  {\begin{tikzpicture}
     \node[classe1=1.6cm] (r) at (0,0) {Rubrica};
     \node[draw=black!80, fill=white, line width=0.7pt, font=\tiny, inner sep=2pt, minimum height=0.55cm,
           anchor=west] (q) at (r.east) {nome: String};
     \node[classe1=1.6cm] (c) at (4.4,0) {Contatto};
     \draw[navig] (q.east) -- node[molt, above, pos=0.85]{0..1} (c);
   \end{tikzpicture}}
  {Il rettangolino attaccato a \code{Rubrica} è il \textbf{qualificatore}: dato un nome, si trova al più un contatto.
   È il disegno di una \code{Map<String, Contatto>} (\vv{mappa}). La molteplicità vale \textbf{per ogni chiave}.}

\voce{naria}{Associazione n-aria}{tre o più classi, con un rombo \oltre}
  {\begin{tikzpicture}
     \node[draw=black!80, fill=white, line width=0.7pt, diamond, minimum size=0.6cm, inner sep=0pt] (d) at (0,0) {};
     \node[classe1=1.5cm] (g) at (-2.2,0) {Giocatore};
     \node[classe1=1.5cm] (s) at (2.2,0) {Squadra};
     \node[classe1=1.5cm] (a) at (0,-1.2) {Anno};
     \draw[assocl] (g) -- (d); \draw[assocl] (d) -- (s); \draw[assocl] (d) -- (a);
     \node[ruolo] at (0,0.5) {Tesseramento};
   \end{tikzpicture}}
  {Un legame fra \textbf{tre} oggetti insieme (un giocatore, in una squadra, in un anno). Il rombo vuoto \textbf{non} è
   un'aggregazione: è il punto in cui si incontrano le linee. In Java diventa una classe con tre riferimenti, come
   una classe associativa.}

\voce{annidata}{Classe annidata}{cerchio con la croce dal lato della classe esterna \oltre}
  {\begin{tikzpicture}
     \node[classe1=1.6cm] (e) at (0,0) {Ordine};
     \node[classe1=1.6cm] (i) at (3.6,0) {Riga};
     \draw[annida] (e.east) -- (i.west);
     \piuAnnida{($(e.east)+(0.13,0)$)}
     \node[ruolo] at (3.6,-0.5) {oppure \code{Ordine::Riga}};
   \end{tikzpicture}}
  {Una classe dichiarata \textbf{dentro} un'altra (\code{class Ordine \{ class Riga \{\} \}}): la linea con il cerchio e la
   croce. In alternativa si scrive il nome completo \code{Ordine::Riga}. Non dice nulla sui campi: se Ordine contiene
   righe, serve anche l'associazione (\vv{annidatacodice}).}

\voce{template}{Classe template}{parametro in un rettangolo tratteggiato \oltre}
  {\begin{tikzpicture}
     \umlclass[2.4cm]{s}{(0,0)}{Scatola}{- contenuto: T}{+ get(): T}
     \node[draw=black!80, dashed, fill=white, font=\tiny, inner sep=2pt] at ($(s.north east)+(-0.1,0.02)$) {T};
     \node[classe1=2.2cm] (l) at (3.6,0) {Scatola<Libro>};
     \draw[dep] (l.west) -- node[stereo, above]{«bind»} (s.east);
   \end{tikzpicture}}
  {Una classe generica \code{class Scatola<T>}: il parametro \code{T} sta nel rettangolo tratteggiato in alto a
   destra. Una versione con un tipo concreto (\code{Scatola<Libro>}) si collega con «bind» (\textit{T}
   $\rightarrow$ \code{Libro}), oppure si scrive solo come tipo di un attributo (\vv{generica}).}

\begin{attenzione}{Due errori di lettura frequenti}
\begin{itemize}[nosep, leftmargin=*]
  \item Il \textbf{triangolino} accanto al nome ($\blacktriangleright$) dice come si legge; la \textbf{freccia} sulla linea
        dice chi conosce chi. Un diagramma concettuale ha spesso il primo e quasi mai la seconda.
  \item La molteplicità scritta vicino a una classe conta le istanze di \textbf{quella} classe collegate a
        \textbf{una} istanza dell'\textbf{altra}. Per leggerla si parte sempre dalla classe opposta.
\end{itemize}
\end{attenzione}
